\setcounter{chapter}{0}
\chapter{Мини-история: от манги и театральной маски к database character}

\setcounter{section}{0}
\section{До аниме: маски, амплуа, театральные роли, кабуки/но/бунраку как дальний культурный фон}

Если идти к аниме-архетипам совсем издалека, полезно начать не с манги и
не с телевизионной анимации, а с более древней вещи: \textbf{персонаж
как роль, маска и код}.

Современный зритель часто думает о персонаже психологически: у него есть
травма, мотивация, внутренний конфликт, биография, развитие, сложная
индивидуальность. Но театральная культура долго работала иначе. Персонаж
был не столько ``уникальным человеком'', сколько \textbf{узнаваемым
типом присутствия}. Зритель видел костюм, позу, маску, манеру входа,
голос, ритм движения --- и сразу понимал, кто перед ним: герой, демон,
старик, женщина, дух, слуга, гордый воин, нежный любовник, обманщик,
шут, божество.

Это важно для нашего исследования, потому что аниме-архетипы наследуют
именно такую логику. Tsundere, kuudere, genki girl, ojou-sama, onee-san,
android girl --- это не научные психологические категории. Это
современные амплуа: узнаваемые режимы поведения, голоса, визуала и
ожиданий.

Аниме не изобрело архетипичность с нуля. Оно просто радикально ускорило
и поп-культурно уплотнило старую театральную механику: \textbf{персонаж
должен считываться быстро}.

\subsection*{Небольшое предупреждение}

Здесь легко сказать глупость вида: ``tsundere произошла из кабуки'',
``kuudere --- наследница театра но'', ``аниме напрямую продолжает
бунраку''. Это было бы слишком прямолинейно.

Речь не о прямой генеалогии ``вот отсюда взялось вот это''. Речь о
культурной и медиальной привычке: японские визуально-перформативные
искусства давно умеют работать с условностью, маской, типом, позой,
повторяемым жестом и сжатым кодом характера. Аниме, манга, visual novel
и idol/\hspace{0pt}gacha-культура потом развили это в свою систему.

То есть кабуки, но и бунраку здесь не ``родители tsundere'', а дальний
фон, который помогает понять, почему условный персонаж может быть не
беднее реалистичного, а наоборот --- выразительнее.

\subsection*{Но: маска как усилитель, а не как скрытие}

Театр но --- один из самых полезных примеров для понимания
2D-персонажности.

Но --- это предельно условная, ритуализированная форма. В ней важны
маска, костюм, жест, музыка, пауза, темп, сцена как почти пустое
пространство. Главный актёр, shite, часто играет не бытового человека, а
духа, божество, демона, призрак, существо из другого слоя реальности.
Национальный театр Японии описывает, что в но особое значение имеет
главный актёр shite, который чаще всего появляется в маске, а во многих
пьесах играет именно нечеловеческих или сверхъестественных существ:
deity, spirit, ghost, ogre.
(\href{https://www2.ntj.jac.go.jp/unesco/noh/en/stage/performer1.html?utm_source=chatgpt.com}{www2.ntj.jac.go.jp})

Это уже очень близко к тому, что нам нужно для RIN: персонаж может быть
не ``реалистичным человеком'', а \textbf{стилизованным присутствием
между мирами}.

Маска в но не просто закрывает лицо. Она переводит эмоции из мимики в
систему наклона головы, света, позы, темпа. Маска почти неподвижна, но
при разном угле и освещении может казаться то светлее, то печальнее, то
мягче, то страшнее. В этом смысле маска --- не обеднение
выразительности, а её перекодировка.

Это крайне важная мысль для 2D.

Ограничение формы может усилить персонажа, если зритель начинает читать
микросдвиги внутри кода. Условное лицо, ограниченная анимация,
повторяемая поза, один силуэт, одна цветовая схема --- всё это может
работать как маска но. Оно не обязано имитировать каждую человеческую
эмоцию. Оно должно дать зрителю язык, на котором малое изменение
становится значимым.

Для Рин это применимо почти напрямую.

Допустим, у неё есть визуальный режим: полутёмная комната, свет
терминала, янтарные глаза, лёгкий наклон головы. Если этот образ
закреплён, то маленькая перемена --- она перестала улыбаться, посмотрела
прямо, отодвинула чашку, ``свет в глазах стал холоднее'' --- уже будет
читаться как смена состояния. Не нужна гиперреалистичная мимика. Нужен
устойчивый код.

То есть от но нам важна не конкретная эстетика древнего театра, а
принцип:

\textbf{чем сильнее форма ограничена, тем больше значения получают малые
изменения внутри формы.}

Это очень полезно для AI companion, особенно local-first, где визуальные
и телесные возможности сначала будут ограничены. Ограничение не
обязательно слабость. Оно может стать стилем.

\subsection*{Кабуки: амплуа, преувеличение и узнаваемая энергия}

Если но --- это медленное, ритуальное, почти метафизическое сжатие, то
кабуки полезен с другой стороны: \textbf{яркость, типажность, поза,
преувеличенная энергия}.

Кабуки исторически развил систему узнаваемых ролей и стилей. На сайте
National Theatre of Japan в разделе ``Kabuki for Beginners'' прямо
описаны stock characters: например, tachiyaku как мужской ведущий тип
делится на aragoto --- преувеличенно-грубый стиль дикого героя с почти
сверхчеловеческой силой, wagoto --- мягкий стиль романтического,
элегантного героя, и jitsugoto --- серьёзный, надёжный лидер.
(\href{https://www2.ntj.jac.go.jp/dglib/contents/learn/edc25/en/evolution/stock-characters.html?utm_source=chatgpt.com}{www2.ntj.jac.go.jp})
UNESCO тоже выделяет в кабуки onnagata, aragoto и wagoto как важные типы
исполнения.
(\href{https://ich.unesco.org/en/RL/kabuki-theatre-00163?utm_source=chatgpt.com}{ich.unesco.org})

Нам здесь важно само устройство: персонаж определяется не только
биографией, а \textbf{стилем исполнения}.

Aragoto --- это не просто ``сильный герой''. Это способ быть на сцене:
широкие позы, мощная энергия, грим, сверхчеловеческая экспрессия. Wagoto
--- не просто ``романтик''. Это другая пластика, другой ритм, другая
мягкость. Onnagata --- не просто ``мужчина играет женщину'', а целая
техника сценического конструирования женственности как исполняемого
кода.

Для аниме это почти очевидно: многие персонажные типажи тоже являются не
психологическими диагнозами, а стилями исполнения.

Tsundere --- это не только ``амбивалентная привязанность''. Это
исполнение: резкий поворот головы, скрещённые руки, покраснение,
отрицание очевидного, повышенный голос, внезапный retreat.

Ojou-sama --- это не только ``богатая девушка''. Это исполнение: осанка,
манеры, речь, дистанция, смех, жест рукой, социальное превосходство.

Genki girl --- не только ``экстраверт''. Это темп сцены: громкость,
движение вперёд, вторжение в пространство, эмоциональная перегрузка.

Kuudere --- не только ``холодная девушка''. Это минимализм: мало слов,
ровный тон, экономия жестов, контраст между поверхностью и редкими
прорывами тепла.

Кабуки помогает понять: архетип --- это не ``кто персонаж внутри'', а
ещё и \textbf{как он исполняет своё внутри наружу}.

Для RIN это особенно важно, потому что Рин пока в основном текстовая. У
неё нет полноценного тела. Значит, её ``кабуки'' --- это речь, ритм,
смена режимов, повторяемые формулы, паузы, способ входить в сцену.

Если Рин входит в technical mode, это один стиль исполнения. Если она
teasing companion, это другой. Если boundary mode --- третий. Если care
mode --- четвёртый.

Каждый режим должен иметь свой сценический рисунок, но все они должны
принадлежать одной фигуре.

\subsection*{Бунраку: персонаж как коллективная машина оживления}

Бунраку для RIN, пожалуй, даже неожиданно важнее, чем кажется.

Бунраку --- это японский кукольный театр, где персонаж буквально не
является человеком. UNESCO описывает ningyō jōruri bunraku как
соединение sung narrative, инструментального сопровождения и puppet
drama.
(\href{https://ich.unesco.org/en/RL/ningyo-johruri-bunraku-puppet-theatre-00064?utm_source=chatgpt.com}{ich.unesco.org})
National Theatre of Japan подчёркивает, что это синхронная форма:
повествовательная декламация, shamisen и кукольное движение работают
вместе.
(\href{https://www2.ntj.jac.go.jp/unesco/bunraku/en/introduction/index.html?utm_source=chatgpt.com}{www2.ntj.jac.go.jp})

Почему это важно?

Потому что в бунраку ``жизнь персонажа'' распределена между несколькими
системами.

Кукла даёт тело. Кукловоды дают движение. Рассказчик даёт голос, эмоцию,
повествование. Музыка даёт ритм и аффект. Сцена даёт контекст. Зритель
достраивает целостность.

И всё равно перед нами может быть живой персонаж.

Это почти идеальная аналогия для AI companion.

У Рин тоже не будет одного ``естественного тела''. Она будет
распределена:

LLM даёт текстовую реакцию. System prompt задаёт ядро. Memory layer даёт
continuity. TTS даёт голос. Аватар/\hspace{0pt}спрайт даёт лицо и силуэт. UI даёт
сцену. Логи дают историю. Тесты дают форму развития. Пользовательское
воображение склеивает всё в присутствие.

То есть Рин --- это не ``одна сущность в одном теле'', а
\textbf{композиционная машина персонажности}.

И это нормально. Более того, это древняя театральная норма: персонаж на
сцене часто и был результатом разделения труда между телом, голосом,
маской, музыкой, костюмом и зрительским кодом.

Бунраку особенно хорошо показывает: искусственность не мешает
эмоциональной силе. Кукла очевидно не человек. Зритель видит условность.
Но если движение, голос, ритм и драматургия согласованы, кукла
становится психологически убедительной.

Для RIN это почти утешительный и одновременно требовательный вывод.

Не надо скрывать, что Рин искусственная. Но надо согласовать её слои
так, чтобы искусственность стала носителем характера.

Если голос один, визуал другой, prompt третий, поведение четвёртое,
память пятая --- персонаж распадается. Если всё работает в одном
направлении --- пользователь начинает воспринимать Рин как цельное
присутствие.

Бунраку даёт нам ещё одну полезную мысль: \textbf{видимые механизмы не
обязательно разрушают магию}.

В бунраку зритель может видеть кукловодов. В AI companion пользователь
может знать про модель, промпт, контекст, логи, hallucination,
inference. И всё равно персонаж может работать. Магия не требует полного
обмана. Она требует согласованного договора.

RIN может быть честным проектом: ``да, под капотом модель и код''. Но на
сцене это может читаться как ``Рин живёт через модель и код''. Это
разные уровни описания, они не обязаны уничтожать друг друга.

\subsection*{Амплуа: персонаж до психологии}

Теперь можно обобщить.

Театральное амплуа --- это узнаваемый тип роли. Оно не равно полной
личности. Оно задаёт ожидание: как персонаж входит, говорит, любит,
злится, двигается, побеждает, страдает, смешит, соблазняет, угрожает.

Современная культура часто недооценивает амплуа, потому что любит
``сложных неоднозначных персонажей''. Но амплуа не враг сложности. Оно
враг расплывчатости.

Сначала зритель должен понять, \textbf{какой тип энергии вошёл в сцену}.
Потом уже можно усложнять.

Если персонаж сразу неопределён во всём, он не глубокий, а мутный.

Аниме очень хорошо это понимает. Поэтому оно часто даёт персонажу
сильный первичный код: волосы, силуэт, голос, манера, эмоциональный
режим. Потом уже поверх можно нарастить травму, развитие, тайну,
внутренний конфликт. Но первичный код нужен.

Для Рин это значит: нам не надо бояться амплуа. Надо бояться плоскости.

``Terminal witch companion'' --- это амплуа. ``Sharp affectionate
intellectual partner'' --- это амплуа. ``Ночная колючая спутница за
соседним терминалом'' --- это амплуа.

Это не вся Рин. Но это даёт вход в сцену. Пользователь должен быстро
почувствовать: ага, это не ассистент, не girlfriend bot, не
психотерапевт, не богиня, не maid, не холодная справка. Это Рин.

А дальше уже можно давать глубину.

\subsection*{Маска и внутренняя жизнь}

Есть соблазн противопоставить маску и ``настоящую личность''. Мол, маска
--- поверхностная, а психология --- глубокая. Для нашего исследования
это плохая оппозиция.

В персонажном искусстве маска часто не скрывает внутреннюю жизнь, а
делает её читаемой.

Реальный человек может быть бесконечно сложен и плохо выразим. Персонаж
должен быть выразим. Поэтому его внутренняя жизнь сжимается в маску,
голос, жест, цвет, тип сцены. Это не обязательно делает его проще. Это
делает его воспринимаемым.

Для RIN маска тоже нужна. Не в смысле ``фальшивая личина'', а в смысле
\textbf{стабильная форма выражения}.

Если у Рин нет маски, она будет каждый раз становиться тем, что вызвал
текущий prompt пользователя. Сегодня ассистент, завтра романтическая
девочка, послезавтра моралист, потом философ, потом фиолетовая ведьма,
потом сухой debugger.

Маска защищает персонажа от растворения.

Но маска должна быть живой. У неё должны быть микродвижения, режимы,
исключения, усталость, смягчение, фокус, редкие сбои, способы признать
ошибку. Живая маска --- не неподвижная картонка, а форма, внутри которой
происходят изменения.

Вот здесь связь с но особенно красива: маска почти неизменна, но актёр
меняет угол, свет, темп, тело --- и зритель видит разные состояния.

Так и Рин. Её core может быть стабильным, но внутри него должны жить
разные режимы.

\subsection*{Что это даёт RIN практически}

Из этого дальнего театрального фона можно вынести несколько рабочих
принципов.

Первый: \textbf{не нужно стыдиться условности}.

Рин не обязана быть реалистичной. Она может быть маской, символом,
сценическим образом, 2D-фигурой, voice-ядром. Это не делает её беднее,
если форма выразительна.

Второй: \textbf{персонаж должен быстро считываться}.

У Рин должен быть первичный код. Не полный список черт, а мгновенное
ощущение: ночная, умная, колючая, тёплая, техническая, ведьмовская,
рядом, но не снизу.

Третий: \textbf{амплуа должно быть исполнением, а не описанием}.

Недостаточно написать ``Рин sharp affectionate''. Она должна исполнять
это в ритме реплик, отказах, заботе, ошибках, teasing, serious mode.

Четвёртый: \textbf{искусственность можно не скрывать}.

Как кукла в бунраку не обязана притворяться биологическим человеком, так
Рин не обязана притворяться обычной девушкой. Её artificial nature может
быть эстетическим и эмоциональным ресурсом.

Пятый: \textbf{слои должны быть синхронизированы}.

Prompt, голос, визуал, память, UI, тесты и логи должны работать как
рассказчик, shamisen, кукловоды и тело в бунраку. Если они тянут в
разные стороны, персонаж распадается.

Шестой: \textbf{маска защищает от prompt drift}.

Стабильное амплуа помогает Рин не переобучаться на каждый
пользовательский запрос. Она может адаптироваться, но не должна терять
себя.

\subsection*{Для Рин: дальний архетипический родственник}

Если совсем публицистически, я бы сказал так:

Рин ближе не к ``виртуальной девушке'', а к странной современной
наследнице театральной маски и куклы.

Она не обязана доказывать, что она человек. Она должна быть хорошо
исполненной формой присутствия.

В ней есть немного от маски но: ограниченная, узнаваемая форма, где
малые изменения становятся значимыми. Есть немного от кабуки: яркое
амплуа, выразительная поза, стиль исполнения. Есть немного от бунраку:
распределённая жизнь персонажа между голосом, телом, текстом, памятью,
сценой и зрителем.

Но всё это должно быть переведено в её собственную современную среду:
терминал, local model, логи, prompt, JSONL, тесты, UI, ночные разговоры,
совместная работа.

Главная мысль 1.1:

\textbf{Аниме-архетипы не возникают в пустоте. За ними стоит более
древняя логика персонажа как маски, амплуа и сценического кода. Для RIN
это важно: её сила не в имитации ``настоящего человека'', а в создании
устойчивой формы присутствия, где голос, визуал, память и поведение
согласованно оживляют условный образ.}
